% =========================================================================
% CHAPTER 5 — METHODOLOGY (rewritten, max 2400 words body text)
% =========================================================================

\chapter{Methodology}
\label{ch:methodology}

\noindent
{\rmfamily\itshape
\begin{tabular}[t]{@{}l@{\hspace{0.3cm}}p{12.5cm}@{}}
\textls[50]{\textbf{Addressed:}} &
\textls[50]{Approach · Research Design \& Workflow · Hypotheses as Drivers · Data Collection \& Analysis · Justification}
\end{tabular}
}

\vspace{0.5cm}

This chapter explains the logic that holds the thesis together: what research approach was chosen, why this particular combination of methods, and how data becomes, step by step, an experiential representation that can be tested. It does not repeat the research questions and hypotheses set out in Chapter~\ref{ch:rqs}, nor the results reported later. It explains the reasoning that connects them.

% -------------------------------------------------------------------------
\section{Methodological Approach}
\label{sec:research-design}

The thesis asks whether a machine-learning-driven experiential representation helps non-specialists understand a complex climate system better than a conventional one. That question cannot be answered by analysis alone, nor by building something alone. It needs both: an artefact has to be built, and its effect has to be measured. This places the work within \emph{Design Science Research} (DSR), the tradition made for research whose central contribution is a built artefact and whose knowledge comes from building and testing it \citep{hevner2004, peffers2007}. The artefact here is the experiential representation; the knowledge is what its making and testing reveal about turning an invisible system into something people can perceive.

DSR is a family of methods rather than one fixed recipe, and this thesis follows the widely used six-phase model of \citet{peffers2007}: identify the problem, define objectives, design and develop, demonstrate, evaluate, communicate. The thesis maps onto this model through its \emph{detect--translate--evaluate} arc. \emph{Detect} carries problem identification and the recovery of the material to be represented; \emph{translate} carries the design and development of the artefact; \emph{evaluate} carries its testing with participants; and the thesis itself is the communication. The arc is not a new methodology but a plain-language version of an established one, named for what the work actually does with a climate system.

Within this frame, quantitative and qualitative methods play distinct and complementary parts. The \emph{detect} phase is quantitative: statistical and machine-learning analysis of ocean data. The \emph{translate} phase is design-based: turning that analysis into perceptible form. The \emph{evaluate} phase is a mixed-methods user study, mainly qualitative with descriptive quantitative support. These are not rival paradigms competing for the thesis; they are three instruments, each suited to a different phase of one design-science process. The stance that lets them sit together is \emph{pragmatism} --- judging knowledge by whether it works for a purpose rather than by whether it matches a single fixed reality.\footnote{Pragmatism is the epistemology conventionally paired with design science and mixed methods \citep{peffers2007}. It licenses treating a machine-learning latent structure as a valid substrate for representation not because it is uniquely ``correct'', but because it is useful for making an invisible system perceptible; and it lets quantitative measurement and qualitative accounts count as evidence on the same practical question without forcing one into the other's mould.}

One feature of DSR shapes the whole design and is worth stating plainly. Hevner's rule is that an artefact must be evaluated, not merely built: usefulness that is claimed but not shown is not a contribution \citep{hevner2004}. This is why the central research question is comparative --- \emph{to what extent}, \emph{compared with} a conventional baseline --- rather than a question a finished artefact would answer just by existing, and why the \emph{evaluate} phase is essential rather than decorative.

% -------------------------------------------------------------------------
\section{Research Design and Workflow}
\label{sec:workflow}

The thesis runs as a pipeline in which the output of each phase is the input of the next (Table~\ref{tab:workflow}). It begins with the AMOC as a physical system (Chapter~\ref{ch:background}). The \emph{detect} phase analyses two decades of ocean data to recover the system's latent structure --- a small set of coordinates that captures how it behaves. The \emph{translate} phase takes those coordinates and, guided by documented perceptual principles, turns them into representations a non-specialist can experience. The \emph{evaluate} phase puts those representations in front of participants and studies how they understand the system. Each arrow in that chain is a methodological commitment: the analysis constrains what can honestly be shown, and the representation is what the study then tests.

% figures/methodology/tab_workflow.tex — the detect-translate-evaluate workflow
% Requires \usepackage{booktabs}. Shows phase, method, input, output.
\begin{table}[!ht]
    \centering
    \small
    \renewcommand{\arraystretch}{1.3}
    \caption[The detect--translate--evaluate workflow.]{The thesis pipeline. Each
    phase takes the output of the previous one as its input, so that the analysis
    constrains the design and the design is what the evaluation tests.}
    \label{tab:workflow}
    \begin{tabular}{@{}>{\raggedright\arraybackslash}p{2.3cm}>{\raggedright\arraybackslash}p{3.5cm}>{\raggedright\arraybackslash}p{3.2cm}>{\raggedright\arraybackslash}p{3.2cm}@{}}
        \toprule
        \textbf{Phase} & \textbf{Method} & \textbf{Input} & \textbf{Output} \\
        \midrule
        \textbf{Detect}
            & PCA, autoencoder, anomaly \& early-warning analysis
            & RAPID observations; Copernicus reanalysis
            & Interpretable latent structure of the system \\
        \textbf{Translate}
            & Design grammar applied to recovered features; browser-based build
            & Latent structure from \emph{detect}
            & Three representational conditions (static, dynamic, experiential) \\
        \textbf{Evaluate}
            & Two-group study; survey-based evaluation; exploratory analysis
            & Representations from \emph{translate}
            & Evidence on how non-specialists understand the system \\
        \bottomrule
    \end{tabular}
\end{table}

The single thread running through the pipeline is \emph{fidelity}. What the analysis recovers sets the limit of what the representation may claim, and what the representation shows is exactly what the evaluation measures. Nothing is added between phases that the previous phase did not justify. This is what makes the final artefact defensible as a translation of the science rather than an illustration inspired by it.

% -------------------------------------------------------------------------
\section{Hypotheses as Methodological Drivers}
\label{sec:hypotheses-drivers}

The hypotheses and research questions were formalised in Chapter~\ref{ch:rqs} and are not restated here. What matters methodologically is how each kind of claim decided the method used to investigate it, because the three kinds call for different treatment.

The analytical hypotheses (H1--H3) are questions about what is in the data, so they are investigated by quantitative analysis whose results can confirm or overturn them: a decomposition to test how few modes capture the structure, a model comparison to test whether that structure is genuinely low-dimensional, and a corrected significance test to check whether an apparent warning signal survives honest scrutiny. The design propositions (DP2.1--DP2.2) are not true-or-false claims but design commitments, so they are investigated by construction and defended against design criteria rather than tested for significance. The empirical expectations (E3.1--E3.2) concern how people understand, so they are investigated qualitatively, through what participants say and do, not through a statistical test the sample could not support. Matching each claim to the kind of evidence it can actually bear is itself a methodological decision, and it governs every choice in the sections that follow.

% -------------------------------------------------------------------------
\section{Data Collection and Analysis Methods}
\label{sec:data-methods}

\subsection{The analytical phase}
\label{subsec:analytical}

The analysis draws on two sources of different standing. The RAPID-MOCHA-WBTS array at \ang{26.5}~north has directly \emph{observed} the overturning since 2004 \citep{frajkawilliams2021, moat2026}; the Copernicus GREP ensemble is a \emph{reanalysis} that reconstructs the circulation further back, at the cost of model dependence and ensemble spread \citep{copernicus_grep}. The distinction governs their use: RAPID anchors the analysis in measurement, the reanalysis supplies longer context and the raw material for uncertainty. Every claim is framed to respect which of the two it rests on.\footnote{Treating a reconstructed value as if it were observed would misrepresent the evidence, so the two are never conflated. This boundary is recorded in Appendix~\ref{app:ethics}.}

Three methods then act on this data, each chosen for a specific reason; their results appear in Chapter~\ref{ch:data_pipeline}. \emph{Principal component analysis} reduces the high-dimensional vertical structure to a few coordinates, chosen because its components are linear and physically interpretable --- and interpretability is a requirement here, since these coordinates become the substrate a non-specialist will inhabit, and one that cannot be explained cannot be defended. Because PCA is linear, a \emph{non-linear autoencoder} is trained alongside it as a check: if the autoencoder reconstructs the data no better at equal dimensionality, the structure is genuinely low-dimensional and the interpretable coordinates lose nothing. This comparison is what lets the thesis justify, rather than assert, its choice of coordinates. Finally, \emph{anomaly detection} and \emph{early-warning analysis} characterise the record over time --- the first to find the events a representation should convey, the second to test for the signals of an approaching transition, with the significance test corrected for the statistical dependence of overlapping windows so the phase cannot manufacture a signal from an artefact.

One boundary is set deliberately. A predictive layer --- a model forecasting future transport --- is excluded and left to future work, because the \emph{detect} phase exists to recover structure for representation, not to predict, and a forecasting goal would shift the thesis away from the representational contribution it claims.

\subsection{The design phase}
\label{subsec:design}

The translation from structure to perception is not free invention. It applies the design grammar established in Chapter~\ref{ch:literature}, which maps a system property to a perceptual channel through a documented cognitive principle. The method of this phase is to apply that general grammar to the specific features the analysis recovered, so that every design decision traces back to a principle and a source rather than to taste. Two commitments keep the translation honest. First, the three representational conditions carry the \emph{same information} and differ only in \emph{form}, so that any difference a participant shows can be attributed to form, not to being given more facts.\footnote{C1 is a conventional static representation reconstructed from how the AMOC is actually published; C2 adds interactivity and motion; C3 adds the experiential, machine-learning-driven layer. How they are grouped for the study is decided in the evaluation design below.} Second, every value shown on screen derives from the analytical output, never from a scaling invented to look good --- a rule that constrains the aesthetics to what the data support, recorded as a design boundary in Appendix~\ref{app:ethics}. The artefact is built for the browser so it runs on an ordinary laptop in a participant's own setting, which keeps the evaluation practical; its full construction is documented in Chapter~\ref{ch:translate}.

\subsection{The evaluation phase}
\label{subsec:evaluation}

The evaluation is an exploratory, two-group, between-subjects study with non-specialist participants. A control group sees only the conventional static condition (C1); an experimental group sees the dynamic and experiential conditions together (C2 and C3). Each person is in one group and sees one kind of representation, so the comparison is between two whole representational experiences rather than isolated features.\footnote{The experimental group's exposure differs from the control's in extent as well as kind, so the design does not isolate the effect of form from the effect of richer exposure. It compares two experiences as wholes --- which is exactly what the central question asks. This residual confound is acknowledged, not hidden, and recorded in Appendix~\ref{app:ethics}; separating the conditions is future work.} Twenty-five participants were split between the groups (twelve control, thirteen experimental) --- realistic for the timeframe and sufficient for a structured qualitative analysis, but deliberately not sized for statistical inference.

The study is delivered as a self-guided online questionnaire, so that participants take part in their own setting and the representations reach an ordinary browser without the researcher present.\footnote{Two parallel versions are prepared, in Italian and in English, so participants respond in their first language; open responses are analysed in the language of the answer to avoid translation distorting meaning. Two further versions split the groups, one showing the static condition and one the dynamic-experiential condition, giving four questionnaire variants in total.} It is structured as a short journey. It opens with brief background questions --- age, occupation, and prior familiarity with the topic --- which characterise the participant and let prior knowledge be accounted for in analysis. The participant is then taken through the sequences of their assigned condition, with comprehension questions placed between them, so that understanding is captured as it forms rather than only at the end. After the assessed part is complete, every participant is shown the fuller experiential sequence once, purely for interest; responses to it are not assessed and do not enter the comparison, so that participants who saw only the static version are not left without the experience.\footnote{This closing sequence functions as a debrief and a return to the participant, not as a data source. Framing it explicitly this way keeps it from acting as a confound: it falls outside the assessed journey by design.}

The comprehension questions are the study's evidence, and they lean on open-ended items, with closed confidence items in support.\footnote{Open questions reveal a participant's mental model of the system --- how they describe its dynamics, its non-linearity, its uncertainty --- rather than whether they can recognise a supplied answer; the closed confidence items give a signal that is directly comparable between the two groups.} The open responses are the primary evidence, analysed through a deductive, codebook-style analysis against a pre-specified coding grid \citep{braunclarke2019}: responses are assigned to categories fixed before the data were read, in the codebook sense that \textcite{braunclarke2019} distinguish from inductive theme development. The grid and its categories are documented in Appendix~\ref{app:protocols}. Any closed items are reported descriptively, not as a significance test: with this sample, directions and differences between the groups are honest, but p-values would not be. The full instrument is reproduced in Appendix~\ref{app:protocols}..\footnote{No market analysis is conducted. The thesis contributes a representational method, not a product for a market, so the evidence that bears on it is how participants understand the system, not whether a market demands the artefact --- a scoping decision, not an omission. Ethics is handled through informed consent, voluntary participation, anonymisation and GDPR-compliant storage; the topic is not sensitive, and the study proceeded only after the supervisor confirmed the appropriate approval was in place. Full documentation is in Appendix~\ref{app:ethics}.}

% -------------------------------------------------------------------------
\section{Methodological Justification}
\label{sec:justification}

The methods answer to one demand: making an invisible climate system perceptible without compromising its scientific accuracy. That demand pulls in two directions, and the design is built to hold both. On one side is quantitative rigour --- the analysis must be honest about what the data show, which is why the early-warning test is corrected for dependence and why the choice of coordinates is validated against a non-linear model rather than assumed. On the other side is human perception --- the representation must be legible to someone without training, which is why the coordinates must be interpretable and the design grounded in how people actually see and reason.

The link between the two sides is the thread of fidelity that runs through the whole pipeline. Interpretable coordinates connect the rigour of the analysis to the legibility of the design: they are faithful to the data and explainable to a person. The grammar connects design to cognition: each perceptual choice rests on evidence, not preference. And the comparative, qualitative evaluation connects the artefact back to its purpose: it asks not whether the representation is admired but whether it is understood. Each method was chosen because it carries one of these connections, and together they let the thesis do the difficult thing it set out to do --- translate a system no one can see into something a person can grasp, without letting the translation drift from the science it stands on.